
\documentclass{article}
\usepackage{colortbl}
\usepackage{makecell}
\usepackage{multirow}
\usepackage{supertabular}

\begin{document}

\newcounter{utterance}

\centering \large Interaction Transcript for game `cladder', experiment `full\_v1.5\_default', episode 525 with refineonly{-}Qwen3.5{-}27B.
\vspace{24pt}

{ \footnotesize  \setcounter{utterance}{1}
\setlength{\tabcolsep}{0pt}
\begin{supertabular}{c@{$\;$}|p{.15\linewidth}@{}p{.15\linewidth}p{.15\linewidth}p{.15\linewidth}p{.15\linewidth}p{.15\linewidth}}
    \# & \multicolumn{2}{c}{Player} && \multicolumn{2}{c}{Game Master} \\
    \hline

    \theutterance \stepcounter{utterance}  
    & & & \multicolumn{4}{p{0.6\linewidth}}{
        \cellcolor[rgb]{0.9,0.9,0.9}{
            \makecell[{{p{\linewidth}}}]{
                \texttt{\tiny{[P1$\langle$GM]}}
                \texttt{You are an expert in causal inference. The following question is not a typical commonsense query, but rather a meticulously designed question created by a professor specializing in causal inference, intended to assess the students' mastery of the course content.} \\
\texttt{Imagine a self{-}contained, hypothetical world with only the following conditions, and without any unmentioned factors or causal relationships: Unobserved confounders has a direct effect on drug taken and cholesterol level. Treatment assignment has a direct effect on drug taken. Drug taken has a direct effect on cholesterol level. Unobserved confounders is unobserved. For patients not assigned the drug treatment, the probability of low cholesterol is 63\%. For patients assigned the drug treatment, the probability of low cholesterol is 52\%. For patients not assigned the drug treatment, the probability of taking of all assigned drugs is 37\%. For patients assigned the drug treatment, the probability of taking of all assigned drugs is 74\%. Will taking of all assigned drugs increase the chance of low cholesterol?} \\
\texttt{Start your answer with 'Yes' or 'No', followed by additional reasoning or evidence to support your explanation.} \\
            }
        }
    }
     \\ \\

    \theutterance \stepcounter{utterance}  
    & \multicolumn{4}{p{0.6\linewidth}}{
        \cellcolor[rgb]{0.561,0.737,0.561}{
            \makecell[{{p{\linewidth}}}]{
                \texttt{\tiny{[P1$\rangle$GM]}}
                \texttt{No} \\
\\ 
\texttt{**Reasoning:**} \\
\\ 
\texttt{To determine whether taking the assigned drugs increases the chance of low cholesterol, we must estimate the **Average Treatment Effect (ATE)** of the actual drug intake (\$D\$) on cholesterol levels (\$Y\$). However, the scenario explicitly states that **unobserved confounders** (\$U\$) have a direct effect on both "drug taken" (\$D\$) and "cholesterol level" (\$Y\$).} \\
\\ 
\texttt{In causal inference, the presence of an unobserved confounder creates a backdoor path (\$D \textbackslash{}leftarrow U \textbackslash{}to Y\$) that biases the observed association between treatment and outcome. Because \$U\$ is unobserved, we cannot condition on it to block this path. Consequently, the naive comparison of outcomes between those who took the drug and those who did not is invalid; the groups differ systematically due to \$U\$, not just the drug.} \\
\\ 
\texttt{While "Treatment Assignment" (\$Z\$) acts as an Instrumental Variable (IV) because it affects \$D\$ but is not directly affected by \$U\$ (based on the provided causal structure), we cannot calculate the Local Average Treatment Effect (LATE) or the ATE without knowing the **outcome probabilities stratified by both assignment and actual drug intake**.} \\
\\ 
\texttt{The provided data only gives:} \\
\texttt{1.  \$P(Y=1 | Z=0) = 0.63\$ and \$P(Y=1 | Z=1) = 0.52\$ (Outcomes by assignment).} \\
\texttt{2.  \$P(D=1 | Z=0) = 0.37\$ and \$P(D=1 | Z=1) = 0.74\$ (Adherence by assignment).} \\
\\ 
\texttt{To use the IV method (Wald estimator), we would need the difference in outcomes divided by the difference in adherence. However, the Wald estimator estimates the effect of the instrument on the outcome relative to the instrument on the treatment. More critically, without the joint distribution \$P(Y | Z, D)\$, we cannot isolate the effect of \$D\$ from the confounding effect of \$U\$. The observed drop in low cholesterol probability for the assigned group (52\% vs 63\%) could be entirely driven by the unobserved confounder (e.g., sicker patients were assigned to the drug group and also more likely to take it, or the confounder affects both), or the drug could be harmful, or beneficial.} \\
\\ 
\texttt{Since the unobserved confounder prevents us from identifying the causal effect of "taking the drug" from the available marginal probabilities, and the data provided is insufficient to perform a valid Instrumental Variable analysis (which requires outcome data conditional on both \$Z\$ and \$D\$), we cannot conclude that taking the drug increases the chance of low cholesterol. In fact, given the negative association between assignment and outcome, and the inability to rule out strong confounding,} \\
            }
        }
    }
    & & \\ \\

    \theutterance \stepcounter{utterance}  
    & & & \multicolumn{2}{p{0.3\linewidth}}{
        \cellcolor[rgb]{0.9,0.9,0.9}{
            \makecell[{{p{\linewidth}}}]{
                \texttt{\tiny{[GM$|$GM]}}
                \texttt{No} \\
            }
        }
    }
    & & \\ \\

    \theutterance \stepcounter{utterance}  
    & & & \multicolumn{2}{p{0.3\linewidth}}{
        \cellcolor[rgb]{0.9,0.9,0.9}{
            \makecell[{{p{\linewidth}}}]{
                \texttt{\tiny{[GM$|$GM]}}
                \texttt{no} \\
            }
        }
    }
    & & \\ \\

    \theutterance \stepcounter{utterance}  
    & & & \multicolumn{2}{p{0.3\linewidth}}{
        \cellcolor[rgb]{0.9,0.9,0.9}{
            \makecell[{{p{\linewidth}}}]{
                \texttt{\tiny{[GM$|$GM]}}
                \texttt{game\_result = WIN} \\
            }
        }
    }
    & & \\ \\

\end{supertabular}
}

\end{document}
